\documentclass[12pt]{article}
\usepackage{amsmath,amssymb}
\usepackage[UTF8]{ctex}
\usepackage[hidelinks]{hyperref}
\title{面向固定翼无人机的风险感知轨迹规划：基于多目标两阶段搜索与飞行包线约束细化\\中文研读版}
\author{原稿作者：Zhiwen Lin}
\begin{document}
\maketitle
\section*{说明}
本文件用于中文研读与汇报。为避免对公式、标签、引用和实验占位内容造成任何无意修改，
完整高水平中文精译、汇报提纲与思路把关见配套文件
\texttt{中文版精译\_汇报与思路把关.md}。
下附原始 \LaTeX{} 源码，便于逐式核对。
\section*{中文题目}
面向固定翼无人机的风险感知轨迹规划：基于多目标两阶段搜索与飞行包线约束细化
\section*{中文摘要}
空间非均匀威胁环境下的固定翼无人机任务规划，不仅需要决定飞行航路，还需要确定如何沿所选航路飞行。
本文提出一种风险感知规划框架，将任务偏好驱动的航路选择与受飞行包线约束的轨迹生成相衔接。
硬禁飞区用于定义不可行空间，连续雷达暴露代理量则按飞行时间积分。
通过三个单目标粗搜索建立共享归一化参考，将飞行时间、威胁暴露与燃油消耗统一到多目标搜索代价中。
粗粒度 A* 首先确定总体航路，随后在三维走廊内以相同目标执行细网格重规划。
三次 B 样条用于生成连续候选轨迹，并通过解析曲率与采样几何检查进行筛选；
前向--后向速度分配进一步满足速度、转弯、垂向速度与切向加速度限制。
最终任务指标依据实际速度产生的飞行时间重新积分。
该框架在简化固定翼飞行包线模型下，为任务偏好、几何细化和最终任务评价建立了可追溯的一致性联系。

\section*{章节中文对照}
1 引言；2 问题建模；3 风险感知两阶段轨迹规划框架；4 数值结果与讨论；5 结论。

\newpage
\section*{原始源码（精确核对用）}
\begin{verbatim}
%%
%% Copyright 2007-2024 Elsevier Ltd
%%
%% This file is part of the 'Elsarticle Bundle'.
%% ---------------------------------------------
%%
%% It may be distributed under the conditions of the LaTeX Project Public
%% License, either version 1.3 of this license or (at your option) any
%% later version.  The latest version of this license is in
%%    http://www.latex-project.org/lppl.txt
%% and version 1.3 or later is part of all distributions of LaTeX
%% version 1999/12/01 or later.
%%
%% The list of all files belonging to the 'Elsarticle Bundle' is
%% given in the file `manifest.txt'.
%%
%% Template article for Elsevier's document class `elsarticle'
%% with numbered style bibliographic references
%% SP 2008/03/01
%% $Id: elsarticle-template-num.tex 249 2024-04-06 10:51:24Z rishi $
%%
\documentclass[preprint,12pt]{elsarticle}

%% Use the option review to obtain double line spacing
%% \documentclass[authoryear,preprint,review,12pt]{elsarticle}

%% Use the options 1p,twocolumn; 3p; 3p,twocolumn; 5p; or 5p,twocolumn
%% for a journal layout:
%% \documentclass[final,1p,times]{elsarticle}
%% \documentclass[final,1p,times,twocolumn]{elsarticle}
%% \documentclass[final,3p,times]{elsarticle}
%% \documentclass[final,3p,times,twocolumn]{elsarticle}
%% \documentclass[final,5p,times]{elsarticle}
%% \documentclass[final,5p,times,twocolumn]{elsarticle}

%% For including figures, graphicx.sty has been loaded in
%% elsarticle.cls. If you prefer to use the old commands
%% please give \usepackage{epsfig}

%% The amssymb package provides various useful mathematical symbols
\usepackage{amssymb}
%% The amsmath package provides various useful equation environments.
\usepackage{amsmath}
\usepackage[hidelinks]{hyperref}
%% The amsthm package provides extended theorem environments
%% \usepackage{amsthm}

%% The lineno packages adds line numbers. Start line numbering with
%% \begin{linenumbers}, end it with \end{linenumbers}. Or switch it on
%% for the whole article with \linenumbers.
%% \usepackage{lineno}

\journal{Aerospace Science and Technology}

\begin{document}

\begin{frontmatter}

%% Title, authors and addresses

%% use the tnoteref command within \title for footnotes;
%% use the tnotetext command for theassociated footnote;
%% use the fnref command within \author or \affiliation for footnotes;
%% use the fntext command for theassociated footnote;
%% use the corref command within \author for corresponding author footnotes;
%% use the cortext command for theassociated footnote;
%% use the ead command for the email address,
%% and the form \ead[url] for the home page:
%% \title{Title\tnoteref{label1}}
%% \tnotetext[label1]{}
%% \author{Name\corref{cor1}\fnref{label2}}
%% \ead{email address}
%% \ead[url]{home page}
%% \fntext[label2]{}
%% \cortext[cor1]{}
%% \affiliation{organization={},
%%             addressline={},
%%             city={},
%%             postcode={},
%%             state={},
%%             country={}}
%% \fntext[label3]{}

\title{Risk-aware trajectory planning for fixed-wing UAVs via multi-objective two-stage search and flight-envelope-constrained refinement}

%% use optional labels to link authors explicitly to addresses:
%% \author[label1,label2]{}
%% \affiliation[label1]{organization={},
%%             addressline={},
%%             city={},
%%             postcode={},
%%             state={},
%%             country={}}
%%
%% \affiliation[label2]{organization={},
%%             addressline={},
%%             city={},
%%             postcode={},
%%             state={},
%%             country={}}

\author{Zhiwen Lin}
\affiliation{organization={},%Department and Organization
            addressline={},
            city={Changsha},
            postcode={410073},
            state={Hunan},
            country={China}
            }

%% Abstract
\begin{abstract}
Fixed-wing mission planning in spatially varying threat environments requires
decisions about where to fly and how to traverse the selected route. A
collision-free path alone does not specify the exposure accumulated during
flight or its compatibility with aircraft motion limits. This paper presents
a risk-aware planning framework that connects mission-dependent route
selection with flight-envelope-constrained trajectory generation. Hard
no-fly regions are excluded from the feasible space, while a continuous
radar-exposure surrogate is integrated over flight time. Time, exposure and
fuel estimates are combined using shared normalization references obtained
from three single-objective coarse searches. A coarse A* search establishes
the route, followed by fine-grid replanning within a three-dimensional
corridor under the same weighted objective. Cubic B-spline candidates are
then evaluated using analytical curvature and sampled geometric checks.
Forward--backward speed assignment enforces the modeled speed, turning,
vertical-rate and tangential-acceleration limits. Mission metrics are
recomputed along the final trajectory using the resulting travel times.
The evaluation protocol combines controlled objective and stage comparisons
with a 72-case preference study in a fixed five-zone environment, separating
full two-stage evaluations from coarse-search sensitivity cases.
%% RESULT SENTENCES: 第四章完成后在这里加入2--3句真实结果，替换上一句中的
%% "evaluation protocol"表述为与已完成实验相符的描述。建议依次写：
%% 1. 相同完整协议下，相对时间基线的暴露变化及对应时间/燃油代价；
%% 2. 一组main的细化/最终约束验证结果；3. 72组中实际观察到的偏好取舍。
%% 使用真实数值、单位和工况，不预设降低百分比，不把72组写成72个场景。
The framework provides a traceable connection between planning preferences,
geometric refinement and final mission evaluation within a reduced
fixed-wing flight-envelope model.
\end{abstract}


%% Graphical abstract: use a simplified version of the framework figure when ready.
%% \begin{graphicalabstract}
%% \end{graphicalabstract}

\begin{highlights}
\item Time, exposure and fuel share fixed references across two search resolutions.
\item Fine-grid replanning resolves a coarse route within a three-dimensional corridor.
\item Spline geometry and assigned speeds undergo explicit flight-envelope checks.
\item Final mission costs are recomputed after geometric refinement and retiming.
\end{highlights}

%% Keywords
\begin{keyword}
Fixed-wing UAV \sep
trajectory planning \sep
risk-aware planning \sep
multi-objective optimization \sep
A* search \sep
B-spline

\end{keyword}

\end{frontmatter}

\section{Introduction}
\label{sec:introduction}

\subsection{Background and motivation}
\label{sec:motivation}

Fixed-wing unmanned aerial vehicle (UAV) mission planning involves coupled
decisions about route geometry, travel time and energy expenditure. In a
threat environment, these decisions also determine the duration and
intensity of exposure. Two collision-free routes can therefore have markedly
different mission costs: a detour may avoid a region of high threat intensity
but require more time and fuel, whereas a shorter route may cross that
region for a limited duration. A useful planner must distinguish absolute
prohibitions from traversable regions whose desirability depends on the
mission preference.

A* provides a general mechanism for searching a graph with assigned edge
costs~\cite{hart1968}. Its capacity to represent soft threats depends on the
chosen graph and cost model. In a binary-occupancy, time-oriented formulation,
obstacle checks establish admissibility and travel-time estimates rank the
remaining routes; different threat intensities within free space do not
affect that ranking unless they are included in the cost. The question
addressed here is how to combine this spatial information with time and
fuel preferences while retaining a clear connection to the trajectory
that will ultimately be evaluated.

Three issues arise in this connection. First, exposure is accumulated over
time, so a spatial threat score alone cannot determine the cost of a
variable-speed trajectory. Second, increasing the resolution of a
three-dimensional search expands its computational requirements, motivating
local refinement of a coarse route. Third, collision-free grid edges can
meet at corners that a fixed-wing aircraft cannot follow directly.
Smoothing changes the geometry, and speed assignment changes the time spent
along it. Consequently, the mission metrics of the final trajectory must
be distinguished from the estimates used to select the discrete path.

\subsection{Related work}
\label{sec:related_work}

Risk-aware and multi-objective planning incorporate more information than
binary occupancy. Guo et al.~\cite{guo2022} combine threat strength and path
length in FC-RRT*, using the resulting cost and flight constraints in tree
expansion and parent selection. Their formulation illustrates that
traversable threats and mission-dependent trade-offs are established
planning concepts. For the present problem, the modeling emphasis is on
time-integrated exposure together with flight time and fuel consumption.
Because subsequent speed assignment can alter exposure even when the route
is unchanged, these quantities also require a common final evaluation
protocol.

Hierarchical planning addresses the cost of resolving large spatial
domains. Liu et al.~\cite{liu2026} couple a multigranularity semantic-risk
representation with coarse corridor guidance and fine-level A* search.
Their urban framework also supports replanning after airspace changes; its
voxel-center output requires subsequent smoothing and platform-specific
dynamic verification. This work establishes a close precedent for
risk-aware corridor refinement. The present study concentrates on
fixed-wing trajectory generation after route selection, carrying shared
time--exposure--fuel references through both search resolutions and
evaluating the refined trajectory under its assigned speed profile.

Fixed-wing trajectory methods explicitly address the relation between
geometry and motion. Li et al.~\cite{li2025} decompose constrained trajectory
optimization into path and speed subproblems and use nonlinearity transfer
and convexification to obtain convergent solution procedures. Zhang et
al.~\cite{zhang2026} use an improved A* path to initialize states and controls
for subsequent pseudospectral trajectory optimization. These studies show
that path--speed decomposition and search-to-trajectory integration already
have substantive precedents. Here, a reduced flight-envelope formulation
is used to connect preference-based graph search with spline refinement
and speed assignment, rather than to solve a full aircraft optimal-control
problem. Parametric spline fitting itself follows established
methods~\cite{dierckx1993}; the application-specific issue is how the fitted
curve is constrained, selected and reevaluated.

\subsection{Research focus and contributions}
\label{sec:contributions}

The problem addressed in this paper is the consistency of the transition
from mission-dependent route selection to a constrained fixed-wing
trajectory. This consistency has three concrete requirements: the two
search resolutions must use the same objective references; the smoothed
curve must be checked against the constraints affected by geometric
modification; and final time, exposure and fuel must be calculated using
the assigned motion rather than inherited from the search. The proposed
framework addresses these requirements through the following contributions.

\begin{enumerate}
\item A mission-preference-driven time--exposure--fuel formulation is
constructed for hard--soft threat environments. Hard no-fly regions define
infeasibility, while a continuous radar-exposure surrogate enters the cost
through its time integral. Three single-objective coarse-search anchors
provide fixed reference scales for weighted route selection and subsequent
reporting.

\item A coarse-to-fine three-dimensional planning procedure applies the
same weighted search objective and reference scales at both resolutions.
The coarse route defines a spatial corridor within which a fresh fine-grid
search resolves local geometry. This procedure makes the role of spatial
refinement explicit without interpreting interpolation as replanning.

\item A constrained path-to-trajectory procedure connects cubic B-spline
refinement, analytical curvature evaluation, sampled feasibility checks
and forward--backward speed assignment. Final mission costs are
reintegrated along the refined trajectory, exposing the changes introduced
by geometric modification and retiming instead of treating discrete-path
statistics as executable-trajectory performance.
\end{enumerate}

The numerical study is organized around controlled comparisons in one
prescribed environment. Full two-stage cases examine the objective choice
and the successive planning stages, while a separate 72-case coarse-search
protocol characterizes preference sensitivity and supplies representative
trajectories. This design separates evidence about the complete planning
chain from evidence about weight-dependent behavior. Feasibility throughout
this paper refers to the specified reduced flight envelope and numerical
checks; it does not imply a full flight-dynamics or tracking validation.

Section~\ref{pr:sec:problem} defines the environment, constraints and mission
objective. Section~\ref{sec:method} presents the planning and evaluation
procedure. Section~\ref{sec:results} specifies the controlled comparisons
and preference analysis, and Section~\ref{sec:conclusions} concludes the paper.
%% 引言不预填实验提升数字。第四章实测完成后，可在本节最后补一句最主要发现。
%% 目前引用支撑的是具体方法关系，不宣称已穷尽相关工作或首次组合这些模块。


% Replace the existing Section 2, including its heading, with this file.
% Requires amsmath and amssymb. All pr: labels are local to this replacement.
\section{Problem formulation}
\label{pr:sec:problem}

\subsection{Three-dimensional mission and threat environment}
\label{pr:sec:environment}
The mission is defined in a three-dimensional Cartesian workspace
$\mathcal{W}\subset\mathbb{R}^{3}$. The UAV position is
$\mathbf{p}=[x,y,z]^{\mathsf T}$, where $x$ and $y$ are horizontal
coordinates and $z$ is altitude. The prescribed mission endpoints are
$\mathbf{p}_{s}$ and $\mathbf{p}_{g}$. A discrete search path is written as
\begin{equation}
 \Gamma_d=(\mathbf{p}_0,\ldots,\mathbf{p}_N),\qquad
 \mathbf{p}_0=\mathbf{p}_{s},\quad \mathbf{p}_N=\mathbf{p}_{g},
 \label{pr:eq:discrete}
\end{equation}
whereas the final time-parameterized trajectory is
$\mathbf{r}(t)\in\mathcal{W}$, $t\in[0,T]$. These representations are
distinguished because collision-free graph connectivity does not establish
continuous-trajectory feasibility. The search, geometric refinement and
speed-assignment procedures are described in Section~3.

Hard exclusion regions and traversable soft-threat regions are modeled
separately. Each hard zone is a vertical cylinder,
\begin{equation}
 \Omega_{F,j}=\left\{\mathbf{p}\in\mathcal{W}:
 (x-x_{F,j})^2+(y-y_{F,j})^2\le R_{F,j}^{2},\quad
 0\le z\le H_{F,j}\right\},
 \label{pr:eq:hardzone}
\end{equation}
where $R_{F,j}$ and $H_{F,j}$ are its horizontal radius and height. With
$\Omega_{\mathrm{hard}}=\bigcup_{j=1}^{N_F}\Omega_{F,j}$, the continuous
collision-avoidance requirement is
\begin{equation}
 \mathbf{r}(t)\notin\Omega_{\mathrm{hard}},\qquad t\in[0,T].
 \label{pr:eq:collision}
\end{equation}

For radar $i$ at $(x_{R,i},y_{R,i},0)$, define the slant range
\begin{equation}
 d_i(\mathbf{p})=\sqrt{(x-x_{R,i})^2+(y-y_{R,i})^2+z^2}.
 \label{pr:eq:slant}
\end{equation}
Its coverage is hemispherical within the nonnegative-altitude workspace.
Let $R_i$ be its outer radius and $R_{\mathrm{in},i}=\eta R_i$, with
$0<\eta<1$. For the unit radar scaling factors used in the experiments,
the local exposure field is
\begin{equation}
 P_{d,i}(\mathbf{p})=
 \begin{cases}
  1, & d_i(\mathbf{p})\le R_{\mathrm{in},i},\\
  \left(\dfrac{R_i-d_i(\mathbf{p})}{R_i-R_{\mathrm{in},i}}\right)^{\nu},
       & R_{\mathrm{in},i}<d_i(\mathbf{p})<R_i,\\
  0, & d_i(\mathbf{p})\ge R_i,
 \end{cases}
 \label{pr:eq:exposurefield}
\end{equation}
where $\nu>0$ controls the decay rate. The experiments use $\eta=0.3$
and $\nu=4$. Overlapping fields are combined by
\begin{equation}
 P_d(\mathbf{p})=\max_{1\le i\le N_R}P_{d,i}(\mathbf{p}),
 \qquad 0\le P_d(\mathbf{p})\le1.
 \label{pr:eq:maxfield}
\end{equation}
This is a prescribed dominant-exposure surrogate, rather than a statistical
fusion formula for the probability of detection by at least one radar.
It does not model radar scan timing, track accumulation or measurement
dependence. Hard zones determine admissibility, while the soft field
contributes to the accumulated exposure objective in
Section~\ref{pr:sec:metrics}.

\subsection{Fixed-wing flight constraints}
\label{pr:sec:constraints}
The model assumes still air and constant aircraft mass. Hence the scalar
path speed $v(t)=\|\dot{\mathbf{r}}(t)\|$ is identified with airspeed.
Feasibility refers to the specified geometric and reduced flight-envelope
constraints; it does not imply validation of a six-degree-of-freedom
aircraft model or a closed-loop tracking controller.

\subsubsection{Mission altitude envelope}
Let $d_s(\mathbf{p})$ and $d_g(\mathbf{p})$ denote the horizontal distances
to the start and goal. For transition distance $D_{\mathrm{tr}}$, define
\begin{equation}
 \alpha_e(\mathbf{p})=\operatorname{clip}
 \left(\frac{d_e(\mathbf{p})}{D_{\mathrm{tr}}},0,1\right),
 \qquad e\in\{s,g\},
 \label{pr:eq:transition}
\end{equation}
and $h_e(\mathbf{p})=z_e+(h_{\mathrm{en}}-z_e)\alpha_e(\mathbf{p})$.
Here $h_{\mathrm{en}}$ is the minimum en-route altitude. The local envelope is
\begin{equation}
 h_{\min}(\mathbf{p})=\min\{h_s(\mathbf{p}),h_g(\mathbf{p}),h_{\mathrm{en}}\},
 \qquad h_{\min}(\mathbf{r}(t))\le z(t)\le h_{\max}.
 \label{pr:eq:altitude}
\end{equation}
This construction retains the endpoint altitudes while imposing the
en-route lower bound outside both transition regions. A sufficient
compatibility condition for the ramp slope is
\begin{equation}
 D_{\mathrm{tr}}\ge
 \max_{e\in\{s,g\}}\frac{\max(h_{\mathrm{en}}-z_e,0)}{\tan\gamma_{\max}}.
 \label{pr:eq:ramp}
\end{equation}
This condition concerns the altitude envelope itself; it does not guarantee
that a feasible route exists in the presence of obstacles. In the present
mission, $z_s=z_g=0$ is an altitude reference for airborne trajectory
segments. Ground roll, take-off dynamics and landing dynamics are not modeled.

\subsubsection{Airspeed constraint}
The aircraft must remain above its stall-speed lower bound. The speed
assignment uses a prescribed cruise ceiling $v_{\mathrm{cr}}$, additionally
bounded by the aircraft maximum speed $v_{\max}$:
\begin{equation}
 v_{\mathrm{stall}}\le v(t)\le v_{\mathrm{cap}},\qquad
 v_{\mathrm{cap}}=\min\{v_{\mathrm{cr}},v_{\max}\}.
 \label{pr:eq:speedcap}
\end{equation}
For the reported configuration, $v_{\mathrm{stall}}=33~\mathrm{m\,s^{-1}}$,
$v_{\mathrm{cr}}=55~\mathrm{m\,s^{-1}}$ and
$v_{\max}=83~\mathrm{m\,s^{-1}}$. The airborne initial condition is
$v(0)=v_{\mathrm{stall}}$. No terminal-speed equality is prescribed;
$v(T)$ follows from the local ceilings and the forward--backward speed
assignment. The current speed assignment seeks the available speed under
these ceilings and does not solve a separate fuel-optimal speed-control problem.

\subsubsection{Turning constraint}
For a coordinated level turn, the bank-limited radius satisfies
\begin{equation}
 R_{\min}(v)=\frac{v^2}{g\tan\phi_{\max}}.
 \label{pr:eq:turnradius}
\end{equation}
The implementation applies this relationship as a curvature-based envelope
test on the three-dimensional path. For a regular spline
$\mathbf{q}(u)$, its spatial curvature is
\begin{equation}
 \kappa(u)=\frac{\|\mathbf{q}'(u)\times\mathbf{q}''(u)\|}
                    {\|\mathbf{q}'(u)\|^3},\qquad
 v_{\mathrm{turn}}(u)=\sqrt{\frac{g\tan\phi_{\max}}{\kappa(u)}}.
 \label{pr:eq:curvature}
\end{equation}
The turning ceiling is interpreted as infinite where $\kappa=0$.
The assigned speed must not exceed both $v_{\mathrm{cap}}$ and
$v_{\mathrm{turn}}$. If $v_{\mathrm{turn}}<v_{\mathrm{stall}}$, the geometry
is rejected; the speed is not reduced below stall to make it acceptable.
Using spatial curvature in this level-turn relationship is a reduced-model
approximation for gently climbing trajectories, not an exact bank-angle
reconstruction for general three-dimensional flight. A constant stall
threshold is used; bank-dependent stall speed and full load-factor dynamics
are outside the present model.

\subsubsection{Climb and vertical-rate constraints}
With horizontal arc length $s_h$, the flight-path angle and vertical rate obey
\begin{equation}
 \gamma=\arctan\left(\frac{\mathrm{d}z}{\mathrm{d}s_h}\right),\qquad
 |\gamma|\le\gamma_{\max},\qquad
 -v_{\mathrm{des},\max}\le\dot z\le v_{\mathrm{asc},\max}.
 \label{pr:eq:vertical}
\end{equation}
The angle bound restricts geometry, while $\dot z=v\sin\gamma$ couples
that geometry to speed. Both are required. The current limits are
$\gamma_{\max}=15^\circ$, $v_{\mathrm{asc},\max}=12~\mathrm{m\,s^{-1}}$
and $v_{\mathrm{des},\max}=15~\mathrm{m\,s^{-1}}$.

\subsubsection{Longitudinal acceleration and numerical feasibility}
The tangential acceleration is bounded by
\begin{equation}
 -a_{\mathrm{dec},\max}\le\dot v(t)\le a_{\mathrm{acc},\max},
 \label{pr:eq:acceleration}
\end{equation}
with $a_{\mathrm{acc},\max}=3~\mathrm{m\,s^{-2}}$ and
$a_{\mathrm{dec},\max}=4~\mathrm{m\,s^{-2}}$ in the experiments.
These bounds are imposed during speed assignment rather than added to the cost.

Search edges are screened for environmental and geometric feasibility.
After spline construction, endpoint retention, terminal horizontal heading,
local deviation, altitude, obstacle clearance and curvature are checked
again, followed by speed-envelope validation. The terminal reference heading
is inferred from the incoming discrete path, rather than supplied as an
independent mission heading. The fine-search corridor restricts graph search;
final spline locality is checked against the input path, and membership in
the coarse corridor is diagnostic by default.

The continuous constraints above specify the intended model. Their
implementation uses finite-resolution sampling and numerical tolerances:
the spline geometry is probed at 4096 parameter values, endpoints are accepted
within $2~\mathrm{m}$, and the terminal-heading tolerance is $12^\circ$.
Therefore, acceptance denotes satisfaction of these numerical tests, not
a formal continuous-domain safety certificate. Numerical settings and
resolution checks must accompany the reported experiments.

\subsection{Time, threat-exposure, and fuel models}
\label{pr:sec:metrics}
The mission metrics are flight time $T$, accumulated exposure $C$, and
fuel consumption $F$. They are evaluated on a common trajectory and with
the same time increments. Let $s\in[0,L]$ denote three-dimensional arc
length along the final curve and let $v(s)>0$ be its assigned path speed.
The associated time law is $\mathrm{d}t=\mathrm{d}s/v(s)$.

\subsubsection{Flight time}
The continuous flight-time functional is
\begin{equation}
 T=\int_0^L\frac{1}{v(s)}\,\mathrm{d}s.
 \label{pr:eq:time}
\end{equation}
For final trajectory samples $\mathbf{q}_i$ with speeds $v_i$, the numerical
evaluator uses the chord length
$\ell_i=\|\mathbf{q}_{i+1}-\mathbf{q}_i\|$ and constant tangential
acceleration on each sampled segment. Consequently,
\begin{equation}
 \Delta t_i=\frac{2\ell_i}{v_i+v_{i+1}},\qquad
 a_i=\frac{v_{i+1}^2-v_i^2}{2\ell_i},\qquad
 T_{\mathrm{num}}=\sum_i\Delta t_i,
 \label{pr:eq:segmenttime}
\end{equation}
for nonzero-length segments. This is exact for the adopted segment motion
model, while the sampled representation approximates the continuous spline.
It avoids replacing the final trajectory speed by one mission-wide constant.

During graph search, the speed profile is not yet available. A straight
edge $e$ of length $\ell_e$ and altitude change $\Delta z_e$ instead uses
\begin{equation}
 \Delta t_e^{\mathrm{search}}=
 \max\left\{\frac{\ell_e}{v_{\mathrm{cr}}},
             \frac{|\Delta z_e|}{v_{z,\max}(e)}\right\},
 \label{pr:eq:searchtime}
\end{equation}
where $v_{z,\max}(e)$ is the ascent or descent limit according to the sign
of $\Delta z_e$; the second term is zero for level flight. Here
$v_{\mathrm{cr}}$ lies within the aircraft speed interval. This edge model
is a search surrogate and does not account for acceleration across graph
vertices. The final metrics are recomputed after geometric refinement and
speed assignment; they need not equal the search-stage estimates.

\subsubsection{Accumulated threat exposure}
The exposure objective is
\begin{equation}
 C=\int_0^T P_d(\mathbf{r}(t))\,\mathrm{d}t
   =\int_0^L\frac{P_d(\mathbf{q}(s))}{v(s)}\,\mathrm{d}s.
 \label{pr:eq:exposure}
\end{equation}
Because $P_d$ is dimensionless and bounded by unity, $C$ has units of
seconds and satisfies $0\le C\le T$. It measures time-weighted exposure,
not mission detection probability, interception probability or expected
loss. The same geometric passage can incur different exposure when the
time spent within the field changes. Thus a length integral of $P_d$ alone
would not represent the objective evaluated here.

\subsubsection{Fuel consumption}
Let $v_h=\sqrt{\dot x^2+\dot y^2}$ and $v_z=\dot z$ be the horizontal and
vertical velocity components. The implemented fuel model uses a
quasi-steady parabolic-drag power surrogate. Air density is approximated by
\begin{equation}
 \rho(z)=\rho_0\exp(-z/H_\rho),\qquad
 \rho_0=1.225~\mathrm{kg\,m^{-3}},\quad H_\rho=8500~\mathrm{m}.
 \label{pr:eq:density}
\end{equation}
This is an exponential atmosphere approximation. It is not the full
piecewise International Standard Atmosphere model. The level-flight
mechanical power surrogate is
\begin{equation}
 P_{\mathrm{lev}}(v_h,z)=
 \frac{1}{2}\rho(z)C_{D0}S v_h^3+
 \frac{2(mg)^2}{\rho(z)v_h S\pi\,\mathrm{AR}\,e_O},
 \label{pr:eq:levelpower}
\end{equation}
where $S$ is wing area, $C_{D0}$ is the zero-lift drag coefficient,
$\mathrm{AR}$ is aspect ratio, and $e_O$ is the Oswald efficiency factor.
The first and second terms represent parasitic and induced power,
respectively. The model uses horizontal speed in these terms; it does not
substitute path speed $v$ for $v_h$ in climbing segments.

For the current configuration, descent energy recovery is zero. With
climb efficiency $\eta_c$ and fuel-to-propulsive efficiency $\eta_p$,
the fuel-equivalent power is
\begin{equation}
 P_f(v_h,v_z,z)=\frac{P_{\mathrm{lev}}(v_h,z)+P_c(v_z)}{\eta_p},
 \qquad
 P_c(v_z)=
 \begin{cases}
  mgv_z/\eta_c, & v_z>\delta_z,\\
  0, & v_z\le\delta_z,
 \end{cases}
 \label{pr:eq:fuelpower}
\end{equation}
where $\delta_z=0.01~\mathrm{m\,s^{-1}}$ is the small vertical-rate
deadband used in evaluation. The routine numerically floors $v_h$ at
$v_{\mathrm{stall}}/2$; this floor is inactive for trajectories satisfying
the stated stall and flight-path-angle bounds. With $H_f$ denoting the fuel
lower heating value,
\begin{equation}
 F=\frac{1}{H_f}\int_0^T
 P_f(v_h(t),v_z(t),z(t))\,\mathrm{d}t,
 \qquad H_f=43\times10^6~\mathrm{J\,kg^{-1}}.
 \label{pr:eq:fuel}
\end{equation}
Thus $F$ is reported in kilograms. Since $P_f$ already includes division
by $\eta_p$, the efficiency must not be applied a second time when
converting energy to fuel mass.

For numerical integration, each sampled segment is subdivided at spatial
fractions $0=\lambda_0<\cdots<\lambda_K=1$. Constant tangential acceleration
gives
\begin{equation}
 v_{ij}=\sqrt{(1-\lambda_j)v_i^2+\lambda_j v_{i+1}^2},\qquad
 \Delta t_{ij}=\frac{2\ell_i(\lambda_{j+1}-\lambda_j)}{v_{ij}+v_{i,j+1}}.
 \label{pr:eq:subtime}
\end{equation}
At subsegment midpoint $\mathbf{q}_{ij}^{m}$, the horizontal and vertical
velocities are calculated from the corresponding displacement divided by
$\Delta t_{ij}$. The three numerical metrics are
\begin{equation}
 \begin{aligned}
 T_{\mathrm{num}}&=\sum_{i,j}\Delta t_{ij},\\
 C_{\mathrm{num}}&=\sum_{i,j}P_d(\mathbf{q}_{ij}^{m})\Delta t_{ij},\\
 F_{\mathrm{num}}&=\frac{1}{H_f}\sum_{i,j}
 P_f(v_{h,ij},v_{z,ij},z_{ij}^{m})\Delta t_{ij}.
 \end{aligned}
 \label{pr:eq:quadrature}
\end{equation}
Spatial midpoint evaluation of the nonlinear field and power terms is a
quadrature approximation. Search edges use the same field and power
functions, with their edge time distributed uniformly over the subsegments.

This fuel model is intended for internally consistent comparisons. It
neglects fuel-induced mass change, bank-dependent induced-drag corrections,
transient engine behavior and an explicit tangential-acceleration power
term. These assumptions limit claims about absolute fuel prediction.

\subsection{Anchor-based normalized mission-oriented objective}
\label{pr:sec:objective}
Time, exposure and fuel have different units and numerical scales. A fixed
reference is therefore constructed from three single-objective searches
on the same coarse graph, with identical mission endpoints, environment,
constraints and search-stage evaluation rules. Denote the resulting paths
by $\Gamma_T$, $\Gamma_C$ and $\Gamma_F$, corresponding to the weight
vectors $(1,0,0)$, $(0,1,0)$ and $(0,0,1)$ in $(T,C,F)$ order. These
are discrete search anchors; they are not independently optimized final
continuous trajectories.

Evaluating all three metrics on each anchor gives the payoff table
\begin{equation}
 A=\begin{bmatrix}
 T^{\mathrm{search}}(\Gamma_T)&C^{\mathrm{search}}(\Gamma_T)&F^{\mathrm{search}}(\Gamma_T)\\
 T^{\mathrm{search}}(\Gamma_C)&C^{\mathrm{search}}(\Gamma_C)&F^{\mathrm{search}}(\Gamma_C)\\
 T^{\mathrm{search}}(\Gamma_F)&C^{\mathrm{search}}(\Gamma_F)&F^{\mathrm{search}}(\Gamma_F)
 \end{bmatrix}.
 \label{pr:eq:payoff}
\end{equation}
For objective $k\in\{T,C,F\}$, let the lower reference be the corresponding
anchor's diagonal entry, and let the upper reference be the column maximum:
\begin{equation}
 b_k=A_{kk},\qquad n_k=\max_{a\in\{T,C,F\}}A_{ak},\qquad
 \Delta_k=n_k-b_k>0.
 \label{pr:eq:anchorscale}
\end{equation}
The upper reference is a payoff-table nadir approximation, not the
maximum over all feasible paths or an exact continuous Pareto nadir.
Degenerate reference ranges must be rejected before normalized planning;
they are not silently replaced by an arbitrary denominator.

For a metric vector $\mathbf{f}=(T,C,F)$, define
\begin{equation}
 \widehat T=\frac{T-b_T}{\Delta_T},\qquad
 \widehat C=\frac{C-b_C}{\Delta_C},\qquad
 \widehat F=\frac{F-b_F}{\Delta_F}.
 \label{pr:eq:normalize}
\end{equation}
The same reference values are retained across weights, coarse and fine
searches, and final-trajectory reporting within a fixed experiment.
The normalized values are not clipped to $[0,1]$: a refined or retimed
trajectory can fall outside the three-anchor reference ranges. This does
not invalidate the normalization, but it precludes interpreting the
reference values as exact bounds on final continuous performance.

Mission preference is expressed through
\begin{equation}
 J=w_T\widehat T+w_C\widehat C+w_F\widehat F,\qquad
 w_T,w_C,w_F\ge0,\quad w_T+w_C+w_F>0.
 \label{pr:eq:objective}
\end{equation}
The coefficients specify relative preferences and need not sum to one.
Multiplication of all weights by the same positive constant preserves
path ordering for a fixed model. During graph search, the additive
contribution of edge $e$ is
\begin{equation}
 c_e=\frac{w_T\Delta T_e}{\Delta_T}
     +\frac{w_C\Delta C_e}{\Delta_C}
     +\frac{w_F\Delta F_e}{\Delta_F}.
 \label{pr:eq:edgeobjective}
\end{equation}
The lower-reference offsets are applied once to the complete path:
\begin{equation}
 J^{\mathrm{search}}(\Gamma_d)=\sum_{e\in\Gamma_d}c_e
 -\left(\frac{w_Tb_T}{\Delta_T}+\frac{w_Cb_C}{\Delta_C}
                         +\frac{w_Fb_F}{\Delta_F}\right).
 \label{pr:eq:offset}
\end{equation}
The bracketed term is path independent. Subtracting the lower references
on every edge would instead introduce a spurious dependence on the number
of edges and would not implement Eq.~\eqref{pr:eq:objective}.

This scalarization directs discrete route selection. Spline candidates
are subsequently ranked with a separate geometric refinement score, and
the final $(T,C,F,J)$ values are recomputed using the assigned speed
profile. The sequential procedure is therefore not a joint global
optimization of route geometry, speed and fuel, nor does it establish
global Pareto optimality of the final trajectories.

The sensitivity experiment fixes $w_C=1$ and evaluates
$w_T\in\{0,0.1,\ldots,1.1\}$ and
$w_F\in\{0,0.1,\ldots,0.5\}$, giving 72 configurations.
The coefficients are swept rather than tuned by an additional optimizer.
For tractability, this sweep uses coarse search followed by the same
types of spline refinement, feasibility checks and final evaluation; it
does not repeat the fine-corridor search for all 72 configurations.
Consequently, it characterizes preference sensitivity of that coarse-search
pipeline. Claims about the full two-stage pipeline require separate
controlled evaluations that include the fine-corridor search.


%% Insert after the COMPLETE Section 2, not inside the fuel model.
\section{Proposed risk-aware two-stage trajectory planning framework}
\label{sec:method}

The framework separates route selection, geometric refinement and speed
assignment while retaining a common mission evaluator. The coarse search
first identifies a route under the normalized time--exposure--fuel objective.
A fine search then resolves its local geometry within a three-dimensional
corridor. Cubic B-splines replace the resulting polyline with differentiable
candidates, which are screened against geometric constraints and ranked
using a dimensionless refinement score. Finally, a forward--backward pass
assigns speeds along the selected curve, and the mission metrics are
recomputed under the resulting motion. This decomposition limits the size
of the search state while making the distinction between search estimates
and final trajectory performance explicit.

%% FIGURE PLACEHOLDER: overall framework; use vector PDF for the final artwork.
%% Show mission/environment -> three raw anchors and shared scales -> coarse A*
%% -> 3-D corridor -> fine A* -> spline candidates/repairs -> geometric checks
%% -> minimum geometric score -> strict speed assignment -> final T, C, F, J.
%% Show rejection exits for search, geometry and speed failures. Add a dashed
%% coarse-to-spline branch labelled "72-case sensitivity protocol".
%% \begin{figure}[t]
%% \centering
%% \includegraphics[width=\linewidth]{figures/framework.pdf}
%% \caption{Planning and evaluation workflow. Mission preferences govern graph
%% search, whereas geometric refinement and speed assignment are subsequent
%% stages with explicit feasibility checks.}
%% \label{fig:framework}
%% \end{figure}

\subsection{Multi-objective risk-aware A* search}
\label{sec:astar}

The A* search~\cite{hart1968} represents position on an anisotropic Cartesian lattice
with horizontal and vertical spacings $\delta_{xy}$ and $\delta_z$.
Candidate displacements are drawn from
\begin{equation}
 \mathcal{D}=\{(a\delta_{xy},b\delta_{xy},c\delta_z):
 a,b,c\in\{-1,0,1\},\ (a,b,c)\ne(0,0,0)\}.
 \label{me:eq:neighbors}
\end{equation}
Purely vertical moves are excluded, and a displacement is retained only
if its climb angle satisfies the prescribed bound. The map boundary,
altitude envelope and hard exclusion regions further restrict the
admissible graph. Each retained edge is checked at ten spatial locations,
including its endpoints; the nine intervening subsegments supply the
midpoint quadrature for exposure and fuel. A graph state contains position
only. Heading, curvature and acceleration feasibility are consequently
resolved after search rather than represented as additional state variables.

For a partial path ending at $\mathbf p$, the accumulated search cost is
$G(\mathbf p)=\sum_e c_e$, with $c_e$ defined in
Eq.~\eqref{pr:eq:edgeobjective}. The A* priority is
\begin{equation}
 f(\mathbf p)=G(\mathbf p)+h(\mathbf p),\qquad
 h(\mathbf p)=\frac{w_T}{\Delta_T}
              \frac{\|\mathbf p-\mathbf p_g\|}{v_{\mathrm{cr}}}.
 \label{me:eq:heuristic}
\end{equation}
No heuristic inflation is used. The heuristic retains only a lower bound
on the time contribution: risk and fuel are nonnegative, and the search
time of an edge is at least its length divided by $v_{\mathrm{cr}}$.
For adjacent states $\mathbf p$ and $\mathbf q$, the triangle inequality
therefore gives
\begin{equation}
 h(\mathbf p)-h(\mathbf q)
 \le \frac{w_T\|\mathbf p-\mathbf q\|}{\Delta_Tv_{\mathrm{cr}}}
 \le c_{(\mathbf p,\mathbf q)}.
 \label{me:eq:consistency}
\end{equation}
Thus, the heuristic is consistent for the fixed graph and edge model.
When $w_T=0$, it becomes zero and the primary search reduces to uniform-cost
expansion. Equal primary priorities are ordered using accumulated geometric
length plus Euclidean distance to the goal, providing a deterministic
geometric preference in flat-cost regions. This secondary criterion does
not add a distance penalty to the mission objective.

The constant lower-reference offset in Eq.~\eqref{pr:eq:offset} is restored
when reporting the complete path objective. It is not included repeatedly
in edge costs. The resulting minimum-cost interpretation concerns the
discretized search problem, subject to numerical tolerances and successful
search termination. It does not extend to the subsequently modified and
retimed continuous trajectory.
%% Cite a verified primary A* reference here when finalizing the bibliography.

\subsection{Coarse-to-fine three-dimensional planning}
\label{sec:coarsefine}

A uniform fine lattice over the entire workspace requires many expansions
in regions that are irrelevant to a particular mission preference. The
two-stage search instead uses a coarse path $\Gamma_c$ to restrict the
domain of a second, finer search. For any polyline $\Gamma$, define
\begin{equation}
 d(\mathbf p,\Gamma)=
 \min_{[\mathbf a,\mathbf b]\in\Gamma}\;
 \min_{\tau\in[0,1]}
 \|\mathbf p-\mathbf a-\tau(\mathbf b-\mathbf a)\|.
 \label{me:eq:distance}
\end{equation}
This distance uses the closest point on a segment, rather than the closest
waypoint, and treats horizontal and vertical displacement in the same
three-dimensional metric. The search corridor is
\begin{equation}
 \mathcal C(\Gamma_c,r_c)=
 \{\mathbf p\in\mathcal W:d(\mathbf p,\Gamma_c)\le r_c\}.
 \label{me:eq:corridor}
\end{equation}
Where grid snapping changes an endpoint, the physical endpoint is connected
to the coarse polyline when constructing the corridor reference. The
endpoints used in the present experiments lie on both lattices.

The fine graph retains nodes inside this corridor and samples candidate
edges to prevent a connection between admissible endpoints from leaving
the corridor unchecked. For the reported settings, corridor probes are
spaced no more than approximately $100~\mathrm{m}$ along an edge, with at
least three probes. Altitude and obstacle screening remain active. Both
search stages use the same objective weights and the same anchor-derived
normalization scales; the fine stage does not recompute its own anchors.

The fine result $\Gamma_f$ is obtained by a new graph search, not by
interpolating $\Gamma_c$. Its finer directional and altitude resolution
permits local route changes, but the corridor limits which alternatives
can be explored. The construction therefore trades search coverage for
computational tractability; it does not guarantee equivalence to an
unrestricted fine-grid search. A failed fine search is reported as a
failure rather than relabeling the coarse path as a fine result.

%% FIGURE PLACEHOLDER: coarse/corridor/fine geometry, preferably two panels.
%% Use the SAME actual case: (a) plan view with coarse and fine paths and
%% corridor; (b) 3-D detail or altitude profile showing vertical refinement.
%% Do not draw the 3-D corridor merely as a fixed-width horizontal strip.
%% \begin{figure}[t]
%% \centering
%% \includegraphics[width=\linewidth]{figures/coarse_corridor_fine.pdf}
%% \caption{Coarse-to-fine route refinement within a three-dimensional
%% distance corridor. The corridor restricts the fine search domain.}
%% \label{fig:coarsefine}
%% \end{figure}

\subsection{Cubic B-spline trajectory refinement}
\label{sec:bspline}

The fine polyline satisfies the edge-level search constraints but retains
direction discontinuities at its vertices. Let
$\{\mathbf P_i\}_{i=0}^{n-1}$ denote its ordered vertices after removal of
consecutive duplicates. A cubic parametric B-spline is fitted as
\begin{equation}
 \mathbf q_\sigma(u)=\sum_{j=0}^{K_\sigma-1}
 B_{j,3}(u)\mathbf c_j,\qquad u\in[0,1],
 \label{me:eq:spline}
\end{equation}
where $B_{j,3}$ are cubic basis functions, $\mathbf c_j$ are control points,
and $\sigma$ controls the fit. Chord-length parameterization assigns
$u_i=\sum_{k=0}^{i-1}\|\mathbf P_{k+1}-\mathbf P_k\|/L_{\mathrm{raw}}$,
where $L_{\mathrm{raw}}$ is the polyline length. The initial FITPACK
smoothing fit~\cite{dierckx1993} uses the weighted residual condition
\begin{equation}
 \sum_{i=0}^{n-1}\omega_i^2
       \|\mathbf P_i-\mathbf q_\sigma(u_i)\|^2\le\sigma,
 \qquad
 \omega_i=
 \begin{cases}
 10^5,&i\in\{0,n-1\},\\
 1,&\text{otherwise}.
 \end{cases}
 \label{me:eq:fit}
\end{equation}
Large endpoint weights promote endpoint retention, which is subsequently
verified with an explicit tolerance. The smoothing condition controls the
initial fit; it is not the final mission objective and is not assumed to
remain an exact residual bound after geometric repair.

To explore the balance between polyline fidelity and curvature reduction,
the smoothing budgets are scaled to the input resolution:
\begin{equation}
 \sigma_j=n(\delta_{xy}\beta_j)^2.
 \label{me:eq:budgets}
\end{equation}
The fixed set of 21 dimensionless factors is
\begin{equation}
 \begin{split}
 \mathcal B=\{&0.02,0.03,0.04,0.05,0.0625,0.075,0.09,\\
              &0.105,0.125,0.15,0.18,0.22,0.25,0.35,\\
              &0.50,0.75,1.0,1.5,2.0,3.0,4.0\}.
 \end{split}
 \label{me:eq:budgetset}
\end{equation}
Every budget is examined before selecting a candidate. Small budgets
retain the lattice geometry more closely; larger budgets allow greater
deformation but must still satisfy the local-deviation and flight-envelope
tests. The input is $\Gamma_f$ for the complete two-stage method and
$\Gamma_c$ for the separate coarse-search sensitivity protocol.

Geometric repair addresses constraint violations introduced by smoothing.
Altitude repair raises deficient sampled altitudes relative to the spatial
altitude floor and refits the corrected points. Climb repair limits the
vertical slope using margins below $\tan\gamma_{\max}$ while retaining
the sampled horizontal coordinates before refitting. These operations
generate revised candidates, each of
which is rechecked against the full geometric test set. Repair is a
feasibility procedure and need not reduce the mission objective.

The terminal heading is inferred from the incoming raw route over a
lookback length of $\max(5\delta_{xy},1000~\mathrm{m})$, using the horizontal
chord from the selected upstream vertex to the goal. If the initial
candidate pass finds no geometrically feasible curve and includes a
heading violation, a second pass aligns the final control-polygon edge
with this reference and reevaluates all smoothing budgets. The adjustment
retains the endpoint coordinates and control-point altitudes. The method
thus preserves the approach direction implied by the route; it does not
enforce an independently specified terminal mission heading.

For regular curves with simple interior knots, cubic B-splines provide
$C^2$ geometric continuity and permit curvature evaluation from analytic
derivatives. This property removes polyline corners, but does not by itself
establish feasibility or imply continuous jerk in the time-parameterized
motion. Those issues require the subsequent geometric and speed checks.
%% FITPACK interface details are recorded in literature_alignment.md.

\subsection{Geometric feasibility validation and candidate selection}
\label{sec:trajectory_validation}

Candidate curves are evaluated at $M=4096$ uniformly spaced parameter
values. Their analytic derivatives give spatial curvature through
Eq.~\eqref{pr:eq:curvature} and the local flight-path angle
\begin{equation}
 \gamma(u)=\arctan\!left(
 \frac{|q'_z(u)|}{\sqrt{q'_x(u)^2+q'_y(u)^2}}\right).
 \label{me:eq:splineangle}
\end{equation}
Finite coordinates and nonzero tangents are required. The geometric
turning-radius threshold is
\begin{equation}
 R_{\mathrm{ref}}=
 \max\left\{R_{\mathrm{floor}},
       \frac{v_{\mathrm{stall}}^2}{g\tan\phi_{\max}}\right\}.
 \label{me:eq:radius}
\end{equation}
Requiring $\kappa\le1/R_{\mathrm{ref}}$ prevents acceptance of a turn
that already exceeds the bank envelope at stall speed. It does not require
the curve to support cruise speed everywhere, because speed can subsequently
decrease in curved sections.

A candidate must also satisfy the altitude envelope, avoid the hard
exclusion regions, retain both endpoints within $2~\mathrm{m}$, and match
the terminal horizontal heading within $12^\circ$. Its maximum distance
from the input polyline must satisfy
\begin{equation}
 d(\mathbf q_\sigma(u),\Gamma_{\mathrm{raw}})\le r_s
 \label{me:eq:locality}
\end{equation}
at the validation samples. Here $r_s$ is a refinement locality limit,
distinct from the coarse-search corridor radius $r_c$. In the default
configuration, membership of the final spline in the coarse corridor is
recorded diagnostically, whereas locality relative to the raw input curve
is enforced. Obstacles and the altitude envelope remain hard requirements.

Among the geometrically admissible candidates, selection balances bending,
curvature variation, deviation and length. Let $s$ denote arc length,
$L_\sigma$ the candidate length, and $\kappa_s=\mathrm d\kappa/\mathrm ds$.
The dimensionless components are
\begin{align}
 Q_b&=\frac{1}{L_\sigma}
       \int_0^{L_\sigma}(R_{\mathrm{ref}}\kappa)^2\,\mathrm ds,
       \label{me:eq:bending}\\
 Q_v&=R_{\mathrm{ref}}^3
       \int_0^{L_\sigma}\kappa_s^2\,\mathrm ds,
       \label{me:eq:variation}\\
 Q_d&=\frac{1}{r_s}
       \sqrt{\frac{1}{M}\sum_{m=0}^{M-1}
       d(\mathbf q_\sigma(u_m),\Gamma_{\mathrm{raw}})^2},
       \qquad u_m=\frac{m}{M-1},
       \label{me:eq:deviation}\\
 Q_l&=\frac{L_\sigma}{L_{\mathrm{raw}}}.
       \label{me:eq:length}
\end{align}
The integrals use the cumulative chord-length coordinate of the dense
samples and trapezoidal quadrature; $\kappa_s$ is estimated numerically
on that coordinate. In contrast, $Q_d$ is the root-mean-square deviation
over uniform parameter samples, not an arc-length-weighted integral.
The curvature-variation term penalizes spatial changes in curvature;
it is not a roll-rate constraint.

The refinement score and selection rule are
\begin{equation}
 J_s=0.4Q_b+0.3Q_v+0.2Q_d+0.1Q_l,\qquad
 \mathbf q^*=\mathop{\arg\min}_{\mathbf q\in\mathcal Q_{\mathrm{geom}}}
 J_s(\mathbf q),
 \label{me:eq:selection}
\end{equation}
where $\mathcal Q_{\mathrm{geom}}$ contains the admitted candidates from
the applicable fitting and repair pass. The four coefficients are fixed
geometric design settings, independent of $(w_T,w_C,w_F)$. Selection
enumerates the candidate set rather than stopping at the first feasible
smoothing budget. It is a finite candidate comparison, not a global
optimization over all spline control points.

If no geometrically admissible candidate exists, refinement fails. A
diagnostic fallback curve is not accepted as a mission trajectory. The
selected curve must next pass strict speed validation; the present
procedure does not rerank alternative candidates following a speed-stage
failure. Final acceptance therefore requires both stages to succeed.
Because geometric tests use finite sampling, acceptance has the numerical
meaning stated in Section~\ref{pr:sec:constraints}.

\subsection{Adaptive speed planning and final mission evaluation}
\label{sec:speed}

The selected spline is resampled at approximately uniform arc-length
spacing. A dense set of 8000 parameter probes supplies the cumulative
length map, which is inverted to obtain the final samples
$\{\mathbf q_i\}_{i=0}^{N}$. Their chord lengths are
$\ell_i=\|\mathbf q_{i+1}-\mathbf q_i\|$, and $\kappa_i$ is evaluated
from the spline derivatives at the corresponding parameters. The target
spacing is $100~\mathrm{m}$ in both experimental protocols.

The vertical-rate ceiling of a nondegenerate segment is
\begin{equation}
 \bar v_{z,i}=\begin{cases}
 v_{\mathrm{asc},\max}\ell_i/|z_{i+1}-z_i|,&z_{i+1}>z_i,\\
 v_{\mathrm{des},\max}\ell_i/|z_{i+1}-z_i|,&z_{i+1}<z_i,\\
 +\infty,&z_{i+1}=z_i.
 \end{cases}
 \label{me:eq:verticalcap}
\end{equation}
This ceiling is assigned to both endpoints of the segment. At node $i$,
the admissible speed ceiling is
\begin{equation}
 \bar v_i=\min\left\{
 v_{\mathrm{cr}},v_{\max},
 \sqrt{\frac{g\tan\phi_{\max}}{\kappa_i}},
 \min_{e\in\mathcal E_i}\bar v_{z,e}\right\},
 \label{me:eq:nodecap}
\end{equation}
where $\mathcal E_i$ comprises adjacent segments and the curvature term is
infinite on a straight segment. A ceiling below $v_{\mathrm{stall}}$
causes rejection, rather than a speed assignment below stall.

With $v_0^+=v_{\mathrm{stall}}$, the forward pass propagates the acceleration
bound:
\begin{equation}
 v_{i+1}^+=\min\left\{\bar v_{i+1},
       \sqrt{(v_i^+)^2+2a_{\mathrm{acc},\max}\ell_i}\right\},
 \qquad i=0,\ldots,N-1.
 \label{me:eq:forward}
\end{equation}
The backward pass anticipates downstream speed reductions:
\begin{equation}
 v_N=v_N^+,\qquad
 v_i=\min\left\{v_i^+,
       \sqrt{v_{i+1}^2+2a_{\mathrm{dec},\max}\ell_i}\right\},
 \qquad i=N-1,\ldots,0.
 \label{me:eq:backward}
\end{equation}
No terminal-speed equality is introduced. Applying a segment's vertical
ceiling at both ends bounds its largest vertical speed under the adopted
constant-tangential-acceleration segment model. The final validation checks
speed, turning, climb angle, vertical rate and tangential acceleration;
failure prevents the trajectory from being reported as feasible.

The accepted speed sequence supplies segment times through
Eq.~\eqref{pr:eq:segmenttime}. Time, exposure and fuel are then evaluated
with the common subsegment quadrature in Eq.~\eqref{pr:eq:quadrature},
including renewed altitude and collision screening, and normalized using
the unchanged anchor scales. This last evaluation is essential because
spline deformation changes the spatial threat field sampled by the
vehicle, while retiming changes both residence time and fuel expenditure.
Consequently, a lower geometric refinement score does not by itself imply
a lower final mission objective. The reported output is the accepted
geometry, its assigned speed profile and the recomputed $(T,C,F,J)$;
search-stage estimates are retained separately for diagnosis and stage comparisons.

%% FIGURE PLACEHOLDER: use measured outputs in Section 4.2, not synthetic data.
%% Suggested panels: local raw/spline geometry; analytic spline curvature
%% with 1/R_ref; speed with stall/local ceilings; climb and vertical-rate
%% profiles with limits. Raw-polyline corner curvature is diagnostic only.


\section{Numerical results and discussion}
\label{sec:results}

The evaluation first separates the effects of mission-cost modeling,
spatial refinement and trajectory generation through controlled comparisons
within the same planning framework. A 72-case experiment then examines
mission-preference sensitivity, from which three representative preferences
are selected to illustrate the observed time--exposure--fuel trade-offs.
Search-stage estimates and final trajectory metrics are distinguished
throughout the analysis.

%% 当前紧凑结构：4.1 设置；4.2 三组对比；4.3 四张敏感性核心图；4.4 三个代表。
%% 4.2 的 A/B/C 分别支撑三个贡献，不再为每项贡献另立 subsection。
%% 执行顺序：验收数据 -> 4.1 -> 4.3 -> 4.4自动选择 -> 4.2必要的main对照。
%% B/C共用同一次main；4.4默认直接用sweep中三组最终缓存，不额外跑三次main。
%% 本节英文是比较协议和分析依据；实际结果段落等待实测数据，不预填提升结论。
%% 图均为 %% 占位，不会尝试加载不存在的文件。已合格的72组无需因改稿重跑。

\subsection{Simulation setup and evaluation protocols}
\label{sec:setup}

The simulation uses a $100\times100~\mathrm{km}$ horizontal workspace
containing five ground-based radar sites and five associated cylindrical
exclusion zones. The radar layout is generated with seed 42 and held fixed
throughout the preference comparison. Each radar has a hemispherical
coverage radius of $15~\mathrm{km}$, while its colocated hard zone has a
horizontal radius of $6~\mathrm{km}$ and a height of $5~\mathrm{km}$.
The power-law exposure model uses an inner-radius ratio of 0.3 and an
exponent of 4, with overlapping fields combined by their maximum.
These settings define a controlled synthetic environment; they are not
calibrated measurements of a particular radar or aircraft system.

The mission endpoints are $(2,2,0)$ and $(98,98,0)~\mathrm{km}$.
The en-route altitude interval is $1$--$10~\mathrm{km}$, with a
$5~\mathrm{km}$ horizontal transition distance at each endpoint.
The zero endpoint altitude is the reference used for the airborne
trajectory model, as described in Section~\ref{pr:sec:constraints}.
The vehicle and energy parameters are listed in Table~\ref{tab:aircraft};
the numerical planning settings are summarized in
Table~\ref{tab:planning}.

\begin{table}[t]
 \centering
 \caption{Parameters of the reduced fixed-wing model.}
 \label{tab:aircraft}
 \small
 \begin{tabular}{ll}
 \hline
 Parameter & Value \\
 \hline
 Constant mass & $1100~\mathrm{kg}$ \\
 Wing reference area & $11.5~\mathrm{m^2}$ \\
 Gravitational acceleration & $9.81~\mathrm{m\,s^{-2}}$ \\
 Stall / cruise / aircraft maximum speed & $33 / 55 / 83~\mathrm{m\,s^{-1}}$ \\
 Maximum ascent / descent rate & $12 / 15~\mathrm{m\,s^{-1}}$ \\
 Tangential acceleration / deceleration bound & $3 / 4~\mathrm{m\,s^{-2}}$ \\
 Maximum bank / flight-path angle & $30^\circ / 15^\circ$ \\
 Structural turning-radius floor & $50~\mathrm{m}$ \\
 Zero-lift drag coefficient / aspect ratio & $0.025 / 18$ \\
 Oswald efficiency & $0.85$ \\
 Climb / overall propulsive efficiency & $0.80 / 0.30$ \\
 Descent energy recovery fraction & $0$ \\
 Fuel lower heating value & $43~\mathrm{MJ\,kg^{-1}}$ \\
 Sea-level density / density scale height & $1.225~\mathrm{kg\,m^{-3}} / 8500~\mathrm{m}$ \\
 \hline
 \end{tabular}
\end{table}

\begin{table}[t]
 \centering
 \caption{Search and trajectory-evaluation settings. The sweep and the
 complete two-stage method use different input resolutions and locality limits.}
 \label{tab:planning}
 \small
 \begin{tabular}{ll}
 \hline
 Parameter & Value \\
 \hline
 Coarse horizontal / vertical spacing & $2000 / 500~\mathrm{m}$ \\
 Fine horizontal / vertical spacing & $200 / 50~\mathrm{m}$ \\
 Three-dimensional fine-search corridor radius & $600~\mathrm{m}$ \\
 Spline locality limit, fine input & $200~\mathrm{m}$ \\
 Spline locality limit, coarse sweep input & $3015.90~\mathrm{m}$ \\
 Search edge probes / quadrature subsegments & $10 / 9$ \\
 A* heuristic inflation factor & $1.0$ \\
 Initial smoothing budgets & $21$ \\
 Spline geometric validation samples & $4096$ \\
 Arc-length inversion probes & $8000$ \\
 Final target sampling distance & $100~\mathrm{m}$ \\
 Endpoint / terminal-heading tolerance & $2~\mathrm{m} / 12^\circ$ \\
 Initial speed & $33~\mathrm{m\,s^{-1}}$ \\
 Terminal speed & No prescribed equality \\
 \hline
 \end{tabular}
\end{table}

Two protocols address different experimental questions. The complete
two-stage protocol comprises coarse search, corridor-restricted fine search,
spline refinement and strict speed validation. It supports the controlled
comparisons in Section~\ref{sec:ablation}. The sensitivity protocol
uses coarse search followed by refinement and speed validation for each
of the 72 weight configurations specified in
Section~\ref{pr:sec:objective}. These results support both the sensitivity
analysis and the representative-preference comparison in
Sections~\ref{sec:sensitivity} and~\ref{sec:representative}. Its locality limit is scaled to the coarse
cell diagonal,
\begin{equation}
 r_s^{\mathrm{sweep}}=1.05\sqrt{2(2000)^2+(500)^2}
                    \simeq3015.90~\mathrm{m},
 \label{ex:eq:sweeptrust}
\end{equation}
whereas the complete method uses $r_s=200~\mathrm{m}$ relative to the fine
polyline. The two protocols share the threat and aircraft models, objective
scales, candidate-ranking rule and final metric definitions. Their
different search domains and locality limits are retained explicitly
when interpreting the results.

Search-stage and final-trajectory quantities are distinguished throughout
the analysis. Flight time $T$, exposure $C$ and fuel $F$ are reported in
seconds, exposure-seconds and kilograms, respectively; path length is
reported in kilometres. The exposure integral measures time-weighted
presence in the prescribed field, not a mission-level detection probability.
Cross-preference comparisons use the physical vector $(T,C,F)$: a value
of $J$ computed with one weight vector cannot directly rank a trajectory
scored with another. For a controlled comparison at fixed weights, the
shared anchor normalization provides the common scalar objective.

%% 4.1 操作与核对：
%% A. 环境来自 radar_environment_5zones.py::generate(seed=42)，正文使用 main.py
%%    实际端点 (2,2,0)/(98,98,0) km，不用环境类内部默认的 (5,5)/(95,95)。
%% B. 从运行缓存/原始环境导出五站坐标；归档精确值，不用截图反读坐标。
%% C. 填真实实验机器 CPU、RAM、操作系统、Python/NumPy/SciPy 版本。
%%    本次写作机器的软件版本不能自动充当先前实验的版本。
%% D. 归档 baseline.pkl 中的三个 anchor 原始 T/C/F 与六个尺度参考值；
%%    三条 anchor 未做最终样条处理，不混进“最终轨迹”比较表。
%% E. 计时应分别报告 anchor（一次性）、粗搜索、细搜索、样条、速度/评价。
%%    不把绘图/Excel 导出时间混进搜索时间，不把缓存加载时间当重新计算时间。
%%    只有实际补做重复测量后才能写重复次数、均值/标准差；同图重复计时
%%    反映计算波动，不是多环境统计证据。未测步骤写未记录，不填 0。
%% F. preflight_check.py 可作运行前检查；程序测试不能代替实验数据验收。
%% G. 建议最小数值核查：对选定轨迹另做更密集几何/积分采样，比较最小净空、
%%    最大曲率、最大角度和 T/C/F 的变化；固定几何时无需重跑 A* 或 72 组。
%%    100 -> 50 m 若重新分配速度，属于“离散化+速度策略”联合变化，须单列；
%%    不把它说成仅积分误差。未完成上述核查前不宣称数值收敛或连续安全认证。

%% FIGURE PLACEHOLDER: environment overview; let LaTeX assign figure numbers.
%% Panels: (a) 3-D hemispheres and hard cylinders with mission endpoints;
%% (b) exposure slices at stated altitudes with shared [0,1] color scale.
%% Hard-zone disk only appears when the slice altitude is <= 5 km.
%% Use physical axis units and distinguish hard zones from traversable exposure.
%% \begin{figure}[t]
%% \centering
%% \includegraphics[width=\linewidth]{figures/threat_environment.pdf}
%% \caption{Synthetic mission environment with five radar fields and colocated
%% hard exclusion zones. Soft exposure remains traversable and contributes
%% to the time-integrated objective.}
%% \label{fig:environment}
%% \end{figure}

\subsection{Overall planning performance and controlled comparisons}
\label{sec:ablation}

Three comparisons address distinct parts of the planning chain. The first
examines the effect of the mission objective, the second compares coarse
and fine raw routes from one run, and the third assesses the transition
from that fine route to a continuous trajectory with an assigned speed
profile. The latter two comparisons reuse intermediate outputs rather
than requiring separate planning experiments. These are controlled objective
and stage comparisons, rather than a complete factorial ablation: each
comparison is interpreted at the level of the factors it actually changes.

%% 本节按 A/B/C 三组写，每组遵循：问题 -> 共同条件 -> 数据 -> 实测现象 -> 解释。
%% 不以一张混合 raw/final 的表计算“总提升”，三组对比的指标口径分别明确。
%% 三锚点属于目标建模基础，不单列第四项创新；无须另加外部算法竞赛章节。

\subsubsection{Time-oriented A* versus proposed planning}
\label{sec:objective_comparison}

The reference is a time-oriented A* planner with weights $(1,0,0)$.
Hard-zone admissibility, altitude and climb constraints are retained,
whereas continuous soft exposure and fuel do not enter its search cost.
The proposed formulation uses the normalized composite cost in
Eq.~\eqref{pr:eq:objective}. Thus, the comparison concerns the information
encoded in the objective, rather than a limitation of A* to binary costs.
The radar exposure field remains traversable for both formulations and
is evaluated using the same model after planning.

Both formulations use the same environment, aircraft parameters, lattice
resolutions and refinement and speed-assignment procedures. Their coarse
raw paths can first be compared on the common coarse graph. Final mission
performance is then compared after each route has completed the two-stage
pipeline and passed its feasibility checks. Each route defines its own
fine corridor, so this final comparison measures the end-to-end effect
of the objective rather than holding the fine search domain fixed.

The principal measures are path length $L$, flight time $T$, accumulated
exposure $C$ and fuel consumption $F$. The analysis concerns their
trade-off, not simultaneous improvement in all four quantities. A lower
exposure accompanied by additional time or fuel would indicate the cost
of the selected mission preference, if observed in the data. A scalar
comparison of $J$ requires both trajectories to be evaluated under one
stated weight vector and the same anchor references.

%% A-1：基线是time-oriented A*，不是“A*只懂撞与不撞”；保留所有共同硬约束。
%%      不把整片软雷达区变成基线禁飞区，否则同时改变了可行空间。
%% A-2：多目标权重预先声明；可用4.4规则选出的平衡权重，再执行完整main。
%%      默认(0.3,1,0.2)只能称默认/示例偏好，不自动等于Balanced。
%% A-3：每次main先归档配置、baseline、coarse/final pkl和日志，避免固定文件名覆盖。
%%      原始coarse的T/C/F是search估计；正文主表比较两个通过验证的final结果。
%% A-4：同时加入C和F时只能归因于复合代价；不单独声称两个项各自带来多少收益。
%% A-5：表列Method、weights、L(km)、T(s)、C(exposure-s)、F(kg)，可附总计算时间。
%%      用100*(new-base)/base报告有符号变化；base=0时给绝对差，不加epsilon造百分数。
%% A-6：图用相同场景叠加轨迹，可做俯视/三维/高度多面板；风险切片标明高度。
%%      C_proposed<C_baseline且T_proposed>T_baseline只是可能的权衡，不能提前写成结果。
%%      失败记录保留失败阶段，不能用零值或raw路径代替成功final结果。
%% \begin{figure}[t]
%% \centering
%% \includegraphics[width=\linewidth]{figures/time_vs_risk_aware.pdf}
%% \caption{Time-oriented and mission-oriented final trajectories obtained
%% under the common two-stage planning and flight-envelope validation protocol.}
%% \label{fig:objective_comparison}
%% \end{figure}

\subsubsection{Coarse versus fine planning}
\label{sec:resolution_comparison}

The coarse raw route and the fine raw route are extracted from the same
main run. Both use identical mission weights and anchor scales. Their
comparison therefore concerns the local geometric changes produced by
fine-grid replanning within the three-dimensional corridor, before any
spline deformation or speed reassignment. Spatial overlays and altitude
profiles identify where the fine route departs from the coarse reference.

Search-stage $L$, $T$, $C$, $F$ and $J$ are reported together with the
stage-specific computation times and expansion counts. These observations
assess the practical effect and cost of local refinement; improvement in
every metric is not presumed. Without an unrestricted fine-grid reference,
the results do not establish a quantified speedup over full-domain fine
search or a globally optimal route topology.

%% B-1：使用同一次main的planned_path_coarse.pkl['path','stats']与
%%      planned_path.pkl['path_raw','stats_raw']，不是后者最终path/stats。
%% B-2：叠加粗路径、600m三维走廊和细原始路径，补局部放大/高度剖面。
%%      复用Section 3的coarse_corridor_fine图即可，不重复生成一张同内容图。
%% B-3：计算细原始点/边到粗折线的三维距离；二维投影不能代替该距离。
%%      最终样条可出粗走廊，因默认只强制相对fine raw的200m局部偏离。
%% B-4：表列Coarse raw/Fine raw的搜索估计和阶段开销；fine_time不是两阶段总耗时。
%%      coarse/fine并非简单插值，搜索图也非必然嵌套，故不能预写fine J必定下降。
%% B-5：可说限制了细搜索域并报告实测开销，不能无全域fine对照就声称快了X倍。

\subsubsection{Fine raw path versus B-spline trajectory}
\label{sec:refinement_comparison}

The fine raw path and its refined trajectory are obtained from the same
run. This comparison addresses the distinction between an admissible
piecewise-linear route and a continuous trajectory that passes the adopted
fixed-wing flight-envelope checks. Local geometry, endpoint and heading
retention, altitude and hard-zone screening are examined together with
minimum turning radius, analytic spline curvature and flight-path angle.
Curvature inferred from three raw vertices is a discretization diagnostic,
not continuous curvature at a polyline corner.

The final speed $v(s)$ and signed tangential acceleration $a(s)$ are
assessed against the stall, local curvature, vertical-rate and acceleration
limits. The initial speed is prescribed, while no terminal-speed equality
is imposed. This evaluation concerns numerical compliance with the reduced
flight-envelope model, not validation of a complete aircraft control system.
Final $T$, $C$ and $F$ are recomputed along the refined geometry and its
assigned speed profile. Their simultaneous reduction is not a requirement:
smoothing and retiming serve feasibility, and the geometric selection
score $J_s$ is distinct from the mission objective $J$.

The relation between geometry and exposure is examined along the final
trajectory. Since $dt=ds/v(s)$, a section with elevated threat intensity
contributes more exposure when its traversal speed is lower. Curvature-
or vertical-rate-induced speed reductions can therefore modify the
mission trade-off after route selection. Aligning the threat intensity,
curvature and speed profiles by arc length helps identify whether such
effects occur in the selected case. This interpretation concerns the
observed trajectory; it does not attribute the entire raw-to-final cost
change to a single processing stage.

%% C-1：同一planned_path.pkl取path_raw、path、smooth_result、curvature_smooth、
%%      speed_profile、accel_profile；无需为此再运行main。
%% C-2：图(a)原始/样条局部转弯，(b)解析曲率与1/R_ref，(c)速度与失速/局部上限，
%%      (d)带符号加速度；角度和垂直率可增加小面板或在约束表给极值与阈值。
%% C-3：重点R_min、kappa、gamma、v(s)、a(s)，再列端点误差<=2m、航向误差<=12deg、
%%      最大局部偏离<=200m、altitude余量、硬区检查。不能仅凭曲线平滑宣布可执行。
%% C-4：加速+3/减速-4m/s^2、上升12/下降15m/s分别检查；垂直率包括段内峰值。
%%      原始折线若画非严格速度曲线，明确标成diagnostic，不能充当可飞行基线。
%% C-5：raw的搜索估计与final的变速积分分列；两者差含几何/速度变化，不宣称纯平滑收益。
%%      报告候选预算、选中sigma和修复状态；没有消融则不宣称选择器已优于其他策略。
%% C-6：对至少一个最终案例加密几何/积分采样复核，无需重跑A*；若重新分配速度，
%%      单列为离散化与速度的联合变化。有限采样通过不等于连续安全证明。
%% C-7：复用现有v(s)/曲率图的横轴，标出高暴露弧段，必要时增加小面板Pd(s)。
%%      检查“弯曲/爬升限速 -> 局部dt增加 -> 暴露积分变化”是否真的出现在该案例。
%%      若高暴露区没有明显限速，也如实报告，不能为串联模块强行编造因果故事。
%%      由当前final缓存重积分即可，不需要额外场景或重跑72组。
%% \begin{figure}[t]
%% \centering
%% \includegraphics[width=\linewidth]{figures/refinement_feasibility.pdf}
%% \caption{Fine-route refinement and flight-envelope validation, including
%% analytic spline curvature and the speed and acceleration profiles.}
%% \label{fig:refinement}
%% \end{figure}

\subsection{Mission-preference sensitivity}
\label{sec:sensitivity}

The preference experiment fixes $w_C=1$ and varies
\begin{equation}
 w_T\in\{0,0.1,\ldots,1.1\},\qquad
 w_F\in\{0,0.1,\ldots,0.5\},
 \label{ex:eq:weights}
\end{equation}
resulting in $12\times6=72$ configurations. Each configuration follows
the coarse-search, spline-refinement and speed-assignment protocol in
Section~\ref{sec:setup}. The same environment, aircraft model, anchor
references and final evaluation rules are retained. In particular,
$(w_T,w_F)=(0,0)$ is an exposure-only objective because $w_C$ remains one.

Three heatmaps summarize the final time, exposure and fuel responses to
the two varied coefficients. A fourth plot presents the $T$--$C$ distribution
with fuel consumption encoded by color. Together, these views describe
the objective trade-offs without comparing scalar scores computed under
different preferences. Plateaus and abrupt changes, where present, should
be examined in relation to route selection and subsequent refinement;
monotonic changes in the final metrics are not assumed. The experiment
characterizes the sampled preferences in the stated environment, rather
than a global Pareto frontier or performance across independent scenarios.

%% 4.3只保留下面四张核心图，不再要求正文增加切片图、高度图或行为分区图。
%% 步骤1【验收】：analysis_data_loader核对72个唯一组合、成功状态、schema、平滑版本、
%% run_signature及当前baseline匹配；Excel“全部数据”与对应case缓存一致。
%% 缺失/失败先说明与处理，不能删行后声称72个有效结果，不混用旧工作簿。
%% 步骤2【三幅热力图】：x=w_T,y=w_F，全部12*6网格，单位T(s)、C(exposure-s)、F(kg)。
%% 同一指标的配套图使用一致色标；T/C/F不同量纲，不强行使用同一数值色标。
%% analyze_weight_sensitivity.py已有合并输出；如需三张独立图，可复用同一数据另行导出，
%% 本次只预留图位，并未修改绘图脚本或生成新图。不要从低分辨率截图拼接论文图。
%% 步骤3【散点】：横轴T(s)、纵轴C(exposure-s)、颜色F(kg)，保留所有配置的信息。
%% 重合点可用频数/标记说明；不将72组当独立随机重复实验，不自动添加回归显著性。
%% 步骤4【写作】：先报实测范围与极值位置；再固定一个权重沿行/列解释变化；
%% 最后用散点说明时间--暴露--燃油的代价。检查路线的分析可在后台做，不新增正文图。
%% 若计算非支配点，按三目标<=且至少一项<定义，并公开容差；不要从二维外观断言。
%% 本节的所有指标都是coarse+后处理final，不是72次完整fine流程。
%% \begin{figure}[t]
%% \centering
%% \includegraphics[width=\linewidth]{figures/time_heatmap.pdf}
%% \caption{Final flight time across the 72 configurations of the coarse-search protocol.}
%% \label{fig:time_heatmap}
%% \end{figure}
%% \begin{figure}[t]
%% \centering
%% \includegraphics[width=\linewidth]{figures/exposure_heatmap.pdf}
%% \caption{Time-integrated threat exposure across the 72 preference configurations.}
%% \label{fig:exposure_heatmap}
%% \end{figure}
%% \begin{figure}[t]
%% \centering
%% \includegraphics[width=\linewidth]{figures/fuel_heatmap.pdf}
%% \caption{Final fuel consumption across the 72 preference configurations.}
%% \label{fig:fuel_heatmap}
%% \end{figure}
%% \begin{figure}[t]
%% \centering
%% \includegraphics[width=\linewidth]{figures/time_exposure_scatter.pdf}
%% \caption{Time--exposure trade-offs in the sampled preference set, with fuel
%% consumption represented by color.}
%% \label{fig:tradeoffs}
%% \end{figure}

\subsection{Representative mission preferences}
\label{sec:representative}

Three selection rules identify risk-dominant, time-dominant and balanced
preferences from the same 72-case dataset. For case index $i\in\mathcal I$,
the first two rules select
\begin{equation}
 i_R\in\mathop{\arg\min}_{i\in\mathcal I} C_i,\qquad
 i_T\in\mathop{\arg\min}_{i\in\mathcal I} T_i.
 \label{ex:eq:extremecases}
\end{equation}
For the balanced rule, each physical objective $f_{ik}$,
$k\in\{T,C,F\}$, is scaled within the evaluated batch:
\begin{equation}
 f'_{ik}=\begin{cases}
 \displaystyle\frac{f_{ik}-\min_{j\in\mathcal I}f_{jk}}
 {\max_{j\in\mathcal I}f_{jk}-\min_{j\in\mathcal I}f_{jk}},
 &\max_j f_{jk}>\min_j f_{jk},\\
 0,&\text{otherwise}.
 \end{cases}
 \label{ex:eq:batchscale}
\end{equation}
The balanced case minimizes the Euclidean distance to the batch ideal point,
\begin{equation}
 D_i=\sqrt{(T'_i)^2+(C'_i)^2+(F'_i)^2},\qquad
 i_B\in\mathop{\arg\min}_{i\in\mathcal I}D_i.
 \label{ex:eq:balancedcase}
\end{equation}
This batch scaling is used only for representative selection and does
not replace the three-anchor normalization in the planner. Equal primary
values are resolved lexicographically by $(T,C,F,w_T,w_F)$ for the balanced
rule, $(T,F,w_T,w_F)$ for the exposure rule, and $(C,F,w_T,w_F)$ for the
time rule. These definitions match the implemented selection procedure.

The selected final trajectories are taken directly from the corresponding
sweep caches. One three-dimensional overlay and one table summarize their
geometry, weights and final $(T,C,F)$ values under that common protocol.
The labels describe selection rules within the sampled dataset, not
independent operating environments or globally optimal mission solutions.
If two rules select the same configuration, their coincidence is reported
rather than forcing distinct cases. The balanced choice expresses equal
importance in batch-scaled objective space; it is neither the minimizer
of $J$ across different weights nor a universally preferred mission setting.

%% 步骤1：运行select_representative_weights.py，取“绕飞/快穿/平衡”对应的三行。
%% 代码另外输出“省油”用于诊断，但本节只呈现Risk/Time/Balanced，不另加第四案例。
%% 使用自动规则，不手挑好看的权重；平台邻居数与指标指纹不参与选择。
%% 步骤2：读取对应sweep_cases_v5_3_1逐工况缓存，以TRAJECTORY_FILE选择输入；
%% 三条曲线与表中T/C/F均用同一coarse+后处理协议，不混入main结果。
%% 不把筛选表中四舍五入的值当重新积分输入；从已验收缓存/原始数据取精度。
%% 步骤3：叠加同一环境中的三条最终三维轨迹，统一视角、坐标、硬区与雷达图例。
%% 重合代表如实说明，可以保留三行但注明相同权重/轨迹，不声称三个独立解。
%% 步骤4：表仅Case,w_T,w_C,w_F,T,C,F，T用秒,C用exposure-s,F用kg。
%% 若坚持使用w_R表头，应在全文统一替换；当前w_C对应代码w_risk，避免符号混用。
%% 步骤5：文字用一段解释三种偏好的实际路线与代价，再以一段说明取舍和适用范围。
%% 不称Risk解一定全程绕飞、Time解一定穿越中心；“绕飞/快穿”是脚本名称而非几何定理。
%% 本节不额外跑main。需要完整两阶段验证的权重用于4.2，并单独归档。
%% 下表 -- 是尚未填入的占位，绝非零值；有真实结果后替换并移除此备注。
\begin{table}[t]
 \centering
 \caption{Representative preferences selected from the 72-case coarse-search
 experiment. Metrics refer to the final refined and retimed trajectories.}
 \label{tab:representatives}
 \small
 \begin{tabular}{lrrrrrr}
 \hline
 Case & $w_T$ & $w_C$ & $w_F$ & $T$ (s) & $C$ (exposure-s) & $F$ (kg) \\
 \hline
 Risk-dominant & -- & -- & -- & -- & -- & -- \\
 Time-dominant & -- & -- & -- & -- & -- & -- \\
 Balanced & -- & -- & -- & -- & -- & -- \\
 \hline
 \end{tabular}
\end{table}
%% \begin{figure}[t]
%% \centering
%% \includegraphics[width=\linewidth]{figures/representative_preferences.pdf}
%% \caption{Final three-dimensional trajectories for the risk-dominant,
%% time-dominant and balanced preferences selected from the coarse-search experiment.}
%% \label{fig:representatives}
%% \end{figure}

%% 第四章收束：围绕三个问题总结实测证据，保留效率/暴露/燃油的实际代价；
%% 不预写全面优越、跨场景鲁棒、真实飞行验证或连续安全保证。摘要/结论最后补数字。


\section{Conclusions}
\label{sec:conclusions}

This paper develops a risk-aware planning framework that connects
mission-dependent graph search with fixed-wing trajectory generation.
Hard no-fly constraints and continuous soft exposure have distinct roles
in the model, while time, exposure and fuel are combined through shared
reference scales. Coarse and corridor-constrained fine searches retain
the same mission objective. Subsequent spline refinement and speed
assignment convert the selected path into a trajectory evaluated against
the prescribed flight envelope. Reintegrating the mission metrics after
these stages preserves the distinction between search estimates and the
cost of the resulting motion.

%% EMPIRICAL CONCLUSIONS: 在此插入第四章支持的一个结果段，建议3--4句：
%% (a) 指明主案例权重和时间基线，填写C的变化及T/F实际代价；
%% (b) 指明同次main的粗细路径变化，不把样条/速度的作用归给fine搜索；
%% (c) 给出最终曲率、爬升角、速度/加速度和硬区检查的实际结果；
%% (d) 从72组总结一条本场景内的权重响应规律，注明其为coarse+后处理协议。
%% 不把“需要证明的问题”写成“已经证明的结论”，也不要求三个代价全部下降。

The scope of the study is a prescribed static environment and a reduced
aircraft model. The exposure integral is a planning surrogate, and the
fuel estimate is evaluated under the stated aerodynamic and efficiency
assumptions. Numerical feasibility checks do not constitute a continuous
collision certificate or a closed-loop flight demonstration. Within this
scope, the framework provides an auditable basis for relating mission
preferences to the geometry, motion limits and evaluated cost of a planned
trajectory. Further work will examine the effect of environmental and
model uncertainty and connect the generated references with higher-fidelity
flight dynamics and trajectory tracking.


%% Verified primary references. Included literally in manuscript_updated.tex.
%% Keep citation keys stable when converting this list to a journal BibTeX file.
\begin{thebibliography}{00}

\bibitem{hart1968}
P. E. Hart, N. J. Nilsson, B. Raphael,
A formal basis for the heuristic determination of minimum cost paths,
IEEE Transactions on Systems Science and Cybernetics 4 (2) (1968) 100--107.
\url{https://doi.org/10.1109/TSSC.1968.300136}.

\bibitem{guo2022}
Y. Guo, X. Liu, X. Liu, Y. Yang, W. Zhang,
FC-RRT*: An improved path planning algorithm for UAV in 3D complex environment,
ISPRS International Journal of Geo-Information 11 (2) (2022) 112.
\url{https://doi.org/10.3390/ijgi11020112}.

\bibitem{liu2026}
H. Liu, L. Zeng, M. Lu, K. Tang, B. Li, X. Zhu,
GeoSOT-H-enabled risk-aware hierarchical path planning and emergency
replanning for urban low-altitude UAV missions,
Drones 10 (8) (2026) 603.
\url{https://doi.org/10.3390/drones10080603}.

\bibitem{li2025}
Y. Li, Y. Wang, H. Kuang, X. Liu,
Reliable and efficient trajectory planning for fixed-wing UAVs via
nonlinearity transfer and convexification,
Aerospace Science and Technology 158 (2025) 109919.
\url{https://doi.org/10.1016/j.ast.2024.109919}.

\bibitem{zhang2026}
Z. Zhang, Y. Chen, Y. Bai, H. Hou, H. Jiang,
An efficient trajectory planner for fixed-wing UAVs via improved
A*-guided initialization,
Applied Sciences 16 (18) (2026) 9051.
\url{https://doi.org/10.3390/app16189051}.

\bibitem{dierckx1993}
P. Dierckx, Curve and Surface Fitting with Splines,
Monographs on Numerical Analysis, Oxford University Press, 1993.

\end{thebibliography}

\end{document}

\end{verbatim}
\end{document}
